\subsection{LINEARITY OF THE INTEGRALS OF A LINEAR DIFFERENTIAL EQUATION}

Solving a linear differential equation (2) is a linear problem (Alg., II, p.~240); the homogeneous linear equation

\[
{\frac {d \mathbf {x}}{d t}} = A (t). \mathbf {x}\tag{4}
\]

is said to be associated with the inhomogeneous equation (2); and one knows (Alg., II, p.~241, prop. 14) that if \(u_{1}\) is an integral of the inhomogeneous equation (2) then every integral of this equation is of the form \(u + u_{1}\) where \(u_{1}\) is a solution of the associated homogeneous equation (4), and conversely. We shall first study in this subsection the integrals of a homogeneous equation (4).

PROPOSITION 1. The set I of integrals of the homogeneous linear equation (4), defined on J, is a vector subspace of the space \(\mathcal{C}(\mathrm{J};\mathrm{E})\) of continuous maps from J into E.

The proof is immediate.

THEOREM 2. For every point \((t_0, \mathbf{x}_0)\) of \(\mathbf{J} \times \mathbf{E}\) let \(\mathbf{u}(t, t_0, \mathbf{x}_0)\) be the integral of the homogeneous equation (4) defined on \(\mathbf{J}\) and equal to \(\mathbf{x}_0\) at \(t_0\) .

\(1^{\circ}\) For every point \(t \in \mathbf{J}\) the map \(\mathbf{x}_0 \mapsto \mathbf{u}(t, t_0, \mathbf{x}_0)\) is a bijective bicontinuous linear map \(C(t, t_0)\) of \(\mathbf{E}\) to itself.

\(2^{\circ}\) The map \(t \mapsto C(t, t_0)\) of \(\mathbf{J}\) into \(\mathcal{L}(\mathbf{E})\) is identical to the integral of the homogeneous linear differential equation

\[
\frac {d U}{d t} = A (t) U\tag{5}
\]

which takes the value \(I\) (the identity map of \(E\) to itself) at the point \(t_0\) .

\(3^{\circ}\) For any points \(s, t, u\) of \(\mathbf{J}\)

\[
C (s, u) = C (s, t) C (t, u), \quad C (s, t) = (C (t, s)) ^ {- 1}.\tag{6}
\]

By prop. 1, \(\mathbf{u}(t,t_0,\mathbf{x}_1) + \mathbf{u}(t,t_0,\mathbf{x}_2)\) (resp. \(\lambda \mathbf{u}(t,t_0,\mathbf{x}_0)\) ) is an integral of (4) and takes the value \(\mathbf{x}_1 + \mathbf{x}_2\) (resp. \(\lambda \mathbf{x}_0\) ) at \(t_0\) , so, by th. 1 of IV, p.~179, is identical to \(\mathbf{u}(t,t_0,\mathbf{x}_1 + \mathbf{x}_2)\) (resp. \(\mathbf{u}(t,t_0,\lambda \mathbf{x}_0)\) ); the map \(\mathbf{x}_0\mapsto \mathbf{u}(t,t_0,\mathbf{x}_0)\) is thus a linear map \(C(t,t_0)\) of E into itself, and one can write \(\mathbf{u}(t,t_0,\mathbf{x}_0) = C(t,t_0).\mathbf{x}_0\) .

Since the map \((X, Y) \mapsto XY\) of \(\mathcal{L}(\mathrm{E}) \times \mathcal{L}(\mathrm{E})\) into \(\mathcal{L}(\mathrm{E})\) is continuous (Gen.~Top., X, p.~298, prop. 8), the map \(t \mapsto A(t)U\) of J into \(\mathcal{L}(\mathrm{E})\) is regulated for all \(U \in \mathcal{L}(\mathrm{E})\) ; further (Gen.~Top., X, p.~296)

\[
\| A (t) X - A (t) Y \| = \| A (t) (X - Y) \| \leqslant \| A (t) \| \| X - Y \|,
\]

so one can apply th. 1 of IV, p.~179 to the homogeneous linear equation (5); let \(V(t)\) be the integral of this equation defined on J and equal to \(I\) at \(t_0\) . One has (I, p.~6, prop. 3)

\[
\frac {d}{d t} \big (V (t) \mathbf {x} _ {0} \big) = \frac {d V (t)}{d t} \mathbf {x} _ {0} = A (t) \big (V (t) \mathbf {x} _ {0} \big)
\]

and for \(t = t_{0}\) we have \(V(t)\mathbf{x}_{0} = I\mathbf{x}_{0} = \mathbf{x}_{0}\) ; by th. 1 of IV, p.~179 one must have \(V(t)\mathbf{x}_{0} = C(t, t_{0})\mathbf{x}_{0}\) for all \(x_{0} \in E\) , that is, \(V(t) = C(t, t_{0})\) ; this proves that \(C(t, t_{0})\) belongs to \(\mathcal{L}(\mathrm{E})\) , in other words, that \(x_{0} \mapsto C(t, t_{0})\mathbf{x}_{0}\) is continuous on E, and that the map \(t \mapsto C(t, t_{0})\) is the integral of (5) which is equal to I at \(t_{0}\) .

Finally, the integral \(s \mapsto C(s, u). \mathbf{x}_0\) of (4) is equal to \(C(t, u). \mathbf{x}_0\) at the point \(t\) , so, by definition,

\[
C (s, u). \mathbf {x} _ {0} = C (s, t) \bigl (C (t, u). \mathbf {x} _ {0} \bigr) = \bigl (C (s, t)   C (t, u) \bigr). \mathbf {x} _ {0}
\]

for any \(x_{0} \in E\) , whence the first relation (6); since \(C(s, s) = I\) one has \(C(s, t)C(t, s) = I\) , for any s and t in J; this proves (Set Theory, II, p.~86, corollary) that \(C(t, t_{0})\) is a bijective map of E onto itself, with inverse map \(C(t_{0}, t)\) . This completes the proof of the theorem.

One says that \(C(t, t_0)\) is the resolvent of equation (2) of IV, p.~178.

COROLLARY 1. The map which to every point \(\mathbf{x}_0 \in E\) associates the continuous function \(t \mapsto C(t, t_0). \mathbf{x}_0\) , defined on J, is an isomorphism of the normed space \(E\) onto the vector space \(\mathcal{I}\) of integrals of (4), endowed with the topology of compact convergence.

It is certainly a bijective linear map of E onto \(\mathcal{I}\) ; now \(C(t, t_0)\) is bounded on a compact set \(\mathrm{K} \subset \mathrm{J}\) , so \(\|C(t, t_0). \mathbf{x}_0\| \leqslant \mathrm{M} \| \mathbf{x}_0\|\) for any \(t \in \mathrm{K}\) and \(\mathbf{x}_0 \in \mathrm{E}\) , which shows that this map is continuous; and since

\[
C (t _ {0}, t _ {0}). \mathbf {x} _ {0} = \mathbf {x} _ {0},
\]

it is clear that the inverse map is also continuous.

COROLLARY 2. The map \((s,t)\mapsto C(s,t)\) of \(\mathrm{J}\times \mathrm{J}\) into \(\mathcal{L}(\mathrm{E})\) is continuous.

By (6) we have \(C(s, t) = C(s, t_0) \left( C(t, t_0) \right)^{-1}\) ; now, the map \((X, Y) \mapsto XY\) of \(\mathcal{L}(\mathrm{E}) \times \mathcal{L}(\mathrm{E})\) into \(\mathcal{L}(\mathrm{E})\) is continuous, as is the map \(X \mapsto X^{-1}\) of the (open) group of invertible elements of \(\mathcal{L}(\mathrm{E})\) onto itself (TG, IX, p.~40, prop. 14).

One may note that the map

\[
t \mapsto C (t _ {0}, t) = (C (t, t _ {0})) ^ {- 1}
\]

admits a derivative equal to \(-(C(t,t_{0}))^{-1}(dC(t,t_{0})/dt)(C(t,t_{0}))^{-1}\) (on the complement of a countable set) (I, p.~8, prop. 4), that is to say (by IV, p.~179, formula (5)) equal to \(-C(t_{0},t)A(t)\) .

COROLLARY 3. Let K be a compact interval contained in J, and let \(k = \sup_{t \in \mathbf{K}} \| A(t) \|\) .

For all \(t\) and \(t_0\) in \(\mathbf{K}\)

\[
\left\| C \left(t, t _ {0}\right) - I \right\| \leqslant e ^ {k \left| t - t _ {0} \right|} - 1.\tag{7}
\]

Indeed, \(\|A(t)\mathbf{x}_{0}\|\leqslant k\|\mathbf{x}_{0}\|\) for all \(t\in K\) ; on K the constant function equal to \(x_{0}\) is thus an approximate integral to within \(k\|\mathbf{x}_{0}\|\) by equation (4) of IV, p.~18; by formula (15) of IV, p.~170, one thus has

\[
\left\| C \left(t, t _ {0}\right) \mathbf {x} _ {0} - \mathbf {x} _ {0} \right\| \leqslant \left\| \mathbf {x} _ {0} \right\| \left(e ^ {k | t - t _ {0} |} - 1\right)
\]

for any t and \(t_{0}\) in K, and \(x_{0}\) in E, which is equivalent to the inequality (7) by the definition of the norm on \(\mathcal{L}(\mathrm{E})\) .

PROPOSITION 2. Let B be a continuous endomorphism of E, independent of t, and commuting with \(A(t)\) for all \(t \in J\) ; then B commutes with \(C(t, t_{0})\) for all t and \(t_{0}\) in J.

Indeed, by (5)

\[
\frac {d}{d t} (B C) = B A C = A B C \quad \text { and } \quad \frac {d}{d t} (C B) = A C B,
\]

so \(\frac{d}{dt}\bigl (BC - CB\bigr) = A(BC - CB)\) ; but \(BC(t_0,t_0) - C(t_0,t_0)B = 0\) , so (IV, p.~179, th. 1) \(BC(t,t_0) - C(t,t_0)B = 0\) for all \(t\in \mathrm{J}\) .

An important instance of prop. 2 is that where E is endowed with the structure of a normed vector space with respect to the field of complex numbers C, and where, for every \(t \in J\) , \(A(t)\) is an endomorphism of E for this vector space structure; this means that \(A(t)\) commutes with the continuous endomorphism \(x \mapsto i x\) of E (for the vector space structure over R); then \(C(t, t_{0})\) commutes with this endomorphism, which means that for any t and \(t_{0}\) in J, the map \(C(t, t_{0})\) is a continuous endomorphism for the normed vector space structure of E over C.

